The analysis demonstrates that motor-imagery conditions can yield different dependence-network barcodes within a person. A01 is a clear example: feet and tongue have later component mergers and longer dominant mu loops than the hands. However, the ordering is not uniform across people. The group results therefore support an exploratory, participant-specific description rather than a general movement biomarker. Larger L1 distances for some movement pairs show differences in estimated curves but do not by themselves establish reproducibility or statistical significance.

Several limitations constrain interpretation:
\begin{enumerate}
\item \textbf{Small sample and limited scope.} Nine participants and one training session per participant provide limited power and no independent session-level replication. A01 was already explored when the project was developed; the subsequent summaries and tests are exploratory.
\item \textbf{Sensor mixing and residual artifacts.} Coherence can reflect common sources, volume conduction, reference effects and eye or muscle activity. Original artifact labels and amplitude review do not replace detailed artifact correction. The amplitude sensitivity check concerns matrices only. No causal or anatomical connections can be inferred from these sensor networks.
\item \textbf{Selection and sample size.} Artifact rejection differs across people and conditions, and retained trials need not be a random sample of all trials. Balancing removes surplus data. Repeated selections overlap heavily, and the smallest class is fixed; narrow ranges therefore do not demonstrate strong external reproducibility. The 49-trial check is a single selection.
\item \textbf{Spectral assumptions.} Four-second trials may contain cue-related transients and changing rhythms. One-second Welch windows and band pooling give an average dependence estimate that suppresses temporal and frequency detail. Pooling complex cross spectra also depends on consistency of phase relationships across trials. No sensitivity analysis of window length, band boundaries or reference was performed.
\item \textbf{Topological information loss.} Standard stable ranks discard birth locations. Finite total persistence compresses a whole curve to one number. The nonsignificant omnibus tests do not exclude differences in curve shape, detailed barcodes or specific channel connections. Long-lived features are not automatically neural signal, and short-lived features are not automatically noise.
\item \textbf{Inference assumptions.} Permutation calibration assumes exchangeability of movement summaries within a participant under the null and independence between participants. Trial rejection and temporal/run structure may weaken a randomization interpretation. Holm adjustment covers the four specified primary scalar tests, not every descriptive plot or post hoc interpretation. No held-out classification experiment was conducted.
\end{enumerate}

Useful next steps are replication in the evaluation sessions, more participants, documented artifact correction and sensitivity to spectral settings. A planned full-curve or barcode comparison could complement finite total persistence. If classification is later attempted, training/test separation must be defined at the participant or session level appropriate to the prediction question; overlapping trial selections must not be treated as independent train/test examples.
